\documentclass[12pt,a4paper]{article}

% --- PAQUETES DE CONFIGURACIÓN BÁSICA ---
\usepackage{babel}
\usepackage[utf8]{inputenc}
\usepackage[T1]{fontenc}
\usepackage[margin=2.2cm]{geometry}
\usepackage{hyperref}
\usepackage{enumitem}
\usepackage{booktabs}
\usepackage{xcolor}
\usepackage{titlesec}
\usepackage{titling}
\usepackage{microtype}
\usepackage{array}
\usepackage{multirow}
\usepackage{listings}
\usepackage{tcolorbox}
\usepackage{float}
\usepackage{lmodern} 
\usepackage{eurosym} 

% --- DEFINICIÓN DE COLORES CORPORATIVOS ---
\definecolor{primary}{RGB}{15,118,110}   
\definecolor{secondary}{RGB}{37,99,235} 
\definecolor{muted}{RGB}{107,114,128}     
\definecolor{bgcode}{RGB}{248,250,252}    
\definecolor{frame}{RGB}{226,232,240}     

% --- ESTILO DE SECCIONES ---
\titleformat{\section}{\Large\bfseries\color{primary}}{\thesection}{0.6em}{}
\titleformat{\subsection}{\large\bfseries\color{secondary}}{\thesubsection}{0.6em}{}
\titleformat{\subsubsection}{\bfseries\color{secondary}}{\thesubsubsection}{0.6em}{}

% --- CONFIGURACIÓN DE LISTINGS ---
\lstdefinelanguage{SQLPlus}[]{SQL}{
  morekeywords={IF,ELSE,END,DECLARE,BEGIN,EXCEPTION,WHEN,LOOP,ELSIF,THEN,RAISE,RETURN,ROLE,ACCOUNT,LOCK,UNLOCK,CASE,WHEN,THEN,ELSE,END,CAST,AS,STR_TO_DATE,IFNULL,COALESCE,ON,DUPLICATE,KEY,UPDATE,TRUNCATE},
  sensitive=false
}

\lstset{
  language=SQLPlus,
  basicstyle=\ttfamily\small,
  keywordstyle=\color{secondary}\bfseries,
  commentstyle=\color{muted}\itshape,
  stringstyle=\color{primary},
  columns=fullflexible,
  backgroundcolor=\color{bgcode},
  frame=single,
  rulecolor=\color{frame},
  breaklines=true,
  tabsize=2,
  showstringspaces=false,
  inputencoding = utf8,
  extendedchars = true,
  literate = 
      {á}{{\'a}}1  {é}{{\'e}}1  {í}{{\'i}}1 {ó}{{\'o}}1  {ú}{{\'u}}1
      {Á}{{\'A}}1  {É}{{\'E}}1  {Í}{{\'I}}1 {Ó}{{\'O}}1  {Ú}{{\'U}}1
      {ñ}{{\~n}}1  {Ñ}{{\~N}}1  {¿}{{?`}}1  {¡}{{!`}}1 {€}{{\euro}}1
}

% --- CONFIGURACIÓN DE CAJAS DE TEXTO ---
\tcbset{
  colback=white,
  colframe=primary,
  coltitle=white,
  colbacktitle=primary,
  fonttitle=\bfseries,
  arc=2mm,
  boxrule=0.6pt
}

\title{\textbf{Práctica Avanzada: Saneamiento y Reestructuración (GHA Analytics)}}
\author{Departamento de Informática}
\date{\today}

\begin{document}

\maketitle

\section*{Introducción}
\textbf{Escenario:} La multinacional \textit{Global Health \& Insurance Analytics} (GHA) ha adquirido una clínica regional cuya base de datos es un desastre técnico. Tu misión es liderar el proceso de limpieza y blindaje del sistema antes de conectarlo a la red central.

\section{Fase 0: Preparación del Entorno}
Ejecuta el script \texttt{00\_setup\_gha.sql}. Contiene datos corruptos, duplicados y múltiples violaciones de integridad que deberás resolver.

\section{Escenario 1: Blindaje Estructural (6 Puntos)}
En esta fase, debes aplicar las restricciones de diseño solicitadas. La corrección será automatizada, por lo que el nombre de las restricciones y los tipos de datos deben ser exactos.

\begin{enumerate}
    \item \textbf{Normalización de Identidad (Pacientes):}
    \begin{itemize}
        \item Elimina duplicados exactos en \texttt{pacientes} (mismo NIF y nombre), manteniendo el ID más bajo.
        \item Asegura que el NIF no tenga espacios y cumpla el formato de 8 números y una letra (\textit{Regex}). Los que no cumplan, deben ser corregidos o eliminados.
        \item Convierte la columna \texttt{nif} en \texttt{UNIQUE} y \texttt{NOT NULL}.
    \end{itemize}
    
    \item \textbf{Consistencia de Colegiados (Médicos):}
    \begin{itemize}
        \item Los números de colegiado deben tener el formato \texttt{COL-XX-YYYY} (donde XX es la provincia y YYYY el número). Estandariza los existentes.
        \item Aplica una restricción \texttt{CHECK} para validar este formato.
    \end{itemize}
    
    \item \textbf{Integridad Referencial:}
    \begin{itemize}
        \item Los médicos con especialidades inexistentes deben asignarse a la especialidad 'Medicina General'.
        \item Añade las \texttt{FOREIGN KEY} correspondientes en \texttt{medicos} y \texttt{visitas}.
    \end{itemize}

    \item \textbf{Normalización y División de Tablas:}
    \begin{itemize}
        \item Extrae la información de seguros de la tabla \texttt{pacientes} a una nueva tabla independiente llamada \texttt{seguros\_pacientes}. 
        \item La nueva tabla debe contener el \texttt{paciente\_id}, el \texttt{num\_poliza} y una nueva columna \texttt{estado\_poliza} con el valor por defecto 'ACTIVA'.
        \item Asegura la integridad referencial entre ambas tablas.
    \end{itemize}

    \item \textbf{Columnas Calculadas y Blindaje:}
    \begin{itemize}
        \item Añade a la tabla \texttt{visitas} una columna llamada \texttt{copago\_estimado} de tipo \texttt{DECIMAL(10,2)}.
        \item El valor de esta columna debe calcularse como el 20\% del importe (tras haber saneado la columna \texttt{importe\_sucio}).
        \item \textbf{Blindaje Final:} Una vez saneados los datos, establece como obligatorios (\texttt{NOT NULL}) los campos \texttt{num\_poliza} (en su nueva tabla) y \texttt{copago\_estimado}. Ningún registro debe quedar sin estos valores.
    \end{itemize}

    \item \textbf{Ingesta de Datos Externos:}
    \begin{itemize}
        \item Los registros contenidos en la tabla de \textit{staging} \texttt{raw\_import\_visitas} deben ser procesados e importados a sus tablas correspondientes (\texttt{pacientes} y \texttt{visitas}).
        \item Debes desglosar el campo \texttt{raw\_data}, limpiar los formatos y asegurar que no se creen duplicados si el paciente ya existe.
    \end{itemize}
\end{enumerate}

\section{Escenario 2: Saneamiento Profundo y Memoria (4 Puntos)}
Termina de arreglar todo lo que consideres necesario basándote en la guía de teoría. Debes entregar una \textbf{Memoria Técnica} resumiendo cada cambio realizado y su justificación.

\begin{tcolorbox}[title=Retos sugeridos para el Escenario 2]
\begin{itemize}
    \item \textbf{Gestión de NULLs:} ¿Qué hacer con los pacientes sin email o sin póliza? Usa \texttt{COALESCE} para informes de facturación.
    \item \textbf{Saneamiento de Fechas:} Convierte las fechas de visita (VARCHAR) a tipo \texttt{DATETIME} real. ¡Cuidado con los formatos mezclados!
    \item \textbf{Limpieza Financiera:} Normaliza los importes (quitar €, \$, EUR) y realiza un \texttt{CAST} a \texttt{DECIMAL}.
    \item \textbf{Optimización:} Identifica consultas lentas basándote en el principio de \textit{SARGability}.
\end{itemize}
\end{tcolorbox}

\section*{Instrucciones de Entrega}
\begin{itemize}
    \item \textbf{Script\_Solucion.sql:} Un único archivo con todas las sentencias ordenadas.
    \item \textbf{Memoria\_Tecnica.pdf:} Documentación de los cambios del Escenario 2.
\end{itemize}

\end{document}
